\chapter{Phân tích và thiết kế}
\label{chap:phantich}

Chương này mô hình hóa các quy trình và tương tác của hệ thống (mục~\ref{sec:phantich}), lựa chọn giải pháp công nghệ (mục~\ref{sec:giaiphapcongnghe}) và trình bày thiết kế chi tiết từ kiến trúc, cơ sở dữ liệu, mô hình AI tới API, giao diện, test case và chính sách xử lý với bên thứ ba (mục~\ref{sec:thietke}), dựa trên các yêu cầu đã xác lập tại Chương~\ref{chap:hethong}.

\section{Phân tích}
\label{sec:phantich}

\subsection{Mô hình hóa quy trình nghiệp vụ}
\label{subsec:mohinhquytrinh}

Bảy quy trình nghiệp vụ mô tả tại mục~\ref{sec:quytrinhnghiepvu} được mô hình hóa bằng lưu đồ. Các ký hiệu dùng chung cho toàn bộ lưu đồ được giải thích tại Hình~\ref{fig:chu-thich-flowchart}.

\begin{figure}[H]
\centering
\begin{tikzpicture}[farm/chart,
    desc/.style={anchor=west, text width=9cm, align=left, font=\farmfont}]
    \node[farm/terminal] (t) at (0,0) {Điểm bắt đầu / kết thúc};
    \node[desc] at (2.8,0) {Trạng thái mở đầu hoặc kết thúc của quy trình.};
    \node[farm/process] (p) at (0,-1.1) {Bước xử lý};
    \node[desc] at (2.8,-1.1) {Một thao tác do tác nhân hoặc hệ thống thực hiện.};
    \node[farm/decision] (d) at (0,-2.5) {Điểm\\quyết định};
    \node[desc] at (2.8,-2.5) {Rẽ nhánh theo điều kiện; nhãn trên mũi tên là kết quả (Có/Không).};
    \node[farm/exception] (e) at (0,-3.9) {Ngoại lệ};
    \node[desc] at (2.8,-3.9) {Tình huống bất thường và cách hệ thống xử lý.};
    \node[farm/connector] (c) at (0,-5.1) {A};
    \node[desc] at (2.8,-5.1) {\textbf{Điểm nối:} hai vòng tròn cùng ký hiệu (A với A, B với B) là hai đầu của cùng một đường nối --- luồng đi ra ở vòng tròn này sẽ đi tiếp tại vòng tròn cùng ký hiệu; dùng khi đường nối quá dài hoặc phải chuyển sang hình khác.};
    \draw[farm/arrow] (-1,-6.4) -- (0.9,-6.4);
    \node[desc] at (2.8,-6.4) {Mũi tên liền: luồng đi chính của quy trình.};
    \draw[farm/arrow,dashed] (-1,-7.1) -- (0.9,-7.1);
    \node[desc] at (2.8,-7.1) {Mũi tên nét đứt: liên hệ tham chiếu, không phải bước kế tiếp.};
\end{tikzpicture}
\caption{Chú giải ký hiệu dùng trong các lưu đồ quy trình nghiệp vụ}
\label{fig:chu-thich-flowchart}
\end{figure}


\clearpage
\subsubsection*{Đăng ký và xác thực nhà bán, kiểm duyệt chứng nhận}
Hình~\ref{fig:flow-dang-ky} minh họa quy trình trên dưới dạng lưu đồ, gồm các điểm kiểm tra OTP, eKYC và kiểm duyệt chứng nhận.
\begin{figure}[H]
\centering
\begin{adjustbox}{max totalsize={\textwidth}{\dimexpr\textheight-4.6cm\relax},center}
\begin{tikzpicture}[farm/chart]
    \node[farm/terminal] (start) at (0,0) {Bắt đầu};
    \node[farm/process] (register) at (0,-1.25) {Nhà bán: đăng ký, khai báo định danh và thông tin trang trại};
    \node[farm/process] (otp) at (0,-2.75) {Hệ thống gửi OTP;\\nhà bán nhập mã xác thực};
    \node[farm/decision] (valid) at (0,-4.5) {OTP hợp lệ?};
    \node[farm/decision] (attempts) at (6,-4.5) {Sai quá 5 lần?};
    \node[farm/exception] (lock) at (6,-6.4) {Hệ thống: khóa tạm thời\\30 phút};
    \node[farm/process] (ekyc) at (0,-6.75) {Nhà bán: xác minh danh tính điện tử (eKYC) --- ảnh giấy tờ định danh và ảnh chân dung};
    \node[farm/process] (evidence) at (0,-8.55) {Nhà bán: tải chứng nhận (nếu có), mã số, cơ quan cấp và hạn dùng};
    \node[farm/process] (review) at (0,-10.4) {Kiểm duyệt: đối chiếu kết quả eKYC và chứng nhận với cơ quan cấp hoặc kiểm tra thủ công};
    \node[farm/decision] (approved) at (0,-12.45) {eKYC và hồ sơ\\đạt yêu cầu?};
    \node[farm/processnarrow] (supplement) at (6,-12.45) {Hệ thống: thông báo lý do của bước không đạt; nhà bán bổ sung};
    \node[farm/process] (verified) at (0,-14.75) {Hệ thống: trạng thái ``Đã xác thực'', kích hoạt quyền đăng bán; gắn nhãn chứng nhận nếu có};
    \node[farm/process] (monitor) at (0,-16.45) {Hệ thống: theo dõi hiệu lực chứng nhận trong quá trình hoạt động};
    \node[farm/processnarrow] (expiry) at (6,-16.45) {Còn dưới 30 ngày: nhắc gia hạn.\\Hết hạn: ẩn nhãn đến khi cập nhật; không gỡ sản phẩm.};
    \node[farm/terminal] (finish) at (0,-18.1) {Hoàn tất đăng ký};
    \node[farm/note] at (2.7,-19.9) {\textbf{Ngoại lệ xuyên suốt:} nếu phát hiện chứng nhận giả mạo qua xác minh hoặc khiếu nại, khóa vĩnh viễn tài khoản, ghi vào danh sách vi phạm và chặn đăng ký lại bằng cùng thông tin định danh.};
    \foreach \source/\target in {start/register,register/otp,otp/valid,ekyc/evidence,evidence/review,review/approved,verified/monitor,monitor/finish}
        \draw[farm/arrow] (\source) -- (\target);
    \draw[farm/arrow] (valid) -- node[farm/label,right] {Có} (ekyc);
    \draw[farm/arrow] (valid) -- node[farm/label,above] {Không} (attempts);
    \draw[farm/arrow] (attempts.north) |- node[farm/label,near start,right] {Không: nhập lại} (otp.east);
    \draw[farm/arrow] (attempts) -- node[farm/label,right] {Có} (lock);
    \draw[line width=0.7pt, draw=black!75] (lock.east) -- (8.6,-6.4) -- node[farm/label,right] {Sau 30 phút} (8.6,-2.75) -- (6,-2.75);
    \draw[farm/arrow] (approved) -- node[farm/label,right] {Có} (verified);
    \draw[farm/arrow] (approved) -- node[farm/label,above] {Không} (supplement);
    \draw[farm/arrow] (supplement.west) -- (3.3,-12.45) -- (3.3,-6.75) -- (ekyc.east);
    \draw[farm/arrow,dashed] (monitor) -- (expiry);
\end{tikzpicture}
\end{adjustbox}
\caption{Lưu đồ quy trình đăng ký, xác thực nhà bán và kiểm duyệt chứng nhận}
\label{fig:flow-dang-ky}
\end{figure}

\clearpage
\subsubsection*{Đăng và kiểm duyệt sản phẩm, sinh mã QR truy xuất}
Hình~\ref{fig:flow-san-pham} minh họa quy trình đăng và kiểm duyệt sản phẩm, bao gồm nhánh hậu kiểm nhanh dành cho nhà bán có uy tín cao.
\clearpage
\begin{figure}[p]
\centering
\begin{tikzpicture}[farm/chart]
    \node[farm/terminal] (start) at (0,0) {Nhà bán đã xác thực};
    \node[farm/process] (create) at (0,-1.5) {Nhà bán: tạo sản phẩm, khai báo giá lẻ/sỉ, sản lượng, ngày thu hoạch và nhật ký theo lô};
    \node[farm/decision] (complete) at (0,-3.6) {Có nhật ký gần nhất\\và ảnh thực tế?};
    \node[farm/side] (supplement) at (6,-3.6) {Nhà bán: bổ sung nhật ký và hình ảnh còn thiếu};
    \node[farm/process] (automatic) at (0,-5.6) {Hệ thống: kiểm tra ảnh trùng lặp, từ khóa cấm và giá bất thường};
    \node[farm/decision] (passed) at (0,-7.4) {Đạt kiểm tra\\sơ bộ?};
    \node[farm/side] (revise) at (6,-7.4) {Chưa duyệt; nhà bán chỉnh sửa thông tin sản phẩm};
    \node[farm/decision] (trusted) at (0,-9.5) {Được áp dụng\\hậu kiểm?};
    \node[farm/side] (manual) at (6,-9.5) {Kiểm duyệt: kiểm tra thủ công sản phẩm mới / nhà bán chưa có lịch sử giao dịch};
    \node[farm/decision] (accepted) at (6,-11.8) {Được duyệt?};
    \node[farm/process] (publish) at (0,-11.8) {Hệ thống: duyệt, công khai sản phẩm; sinh QR riêng cho lô, liên kết hồ sơ, chứng nhận, nhật ký và vận chuyển};
    \node[farm/decision] (changed) at (0,-14) {Thay đổi\\trọng yếu?};
    \node[farm/side] (pending) at (6,-14) {Giá thay đổi vượt 20\% hoặc đổi mô tả chất lượng:\\chuyển chờ duyệt lại};
    \node[farm/terminal] (finish) at (0,-16) {Tiếp tục hiển thị};
    \node[farm/connector] (reviseout) at (9,-7.4) {A};
    \node[farm/connector] (revisein) at (6,-0.2) {A};
    \node[farm/connector] (reviewout) at (9,-14) {B};
    \node[farm/connector] (reviewin) at (6,-5.6) {B};
    \node[farm/note] at (2.7,-17.6) {\textbf{Cập nhật nhỏ:} thay đổi còn hàng/hết hàng không cần duyệt lại. Nhánh hậu kiểm chỉ áp dụng cho nhà bán uy tín theo cơ chế được cho phép; việc kiểm tra được thực hiện sau khi công khai sản phẩm.\\[2pt]\textbf{Điểm nối:} các vòng tròn cùng ký hiệu A hoặc B nối tiếp cùng một luồng.};
    \foreach \source/\target in {start/create,create/complete,automatic/passed,manual/accepted,publish/changed}
        \draw[farm/arrow] (\source) -- (\target);
    \draw[farm/arrow] (complete) -- node[farm/label,right] {Có} (automatic);
    \draw[farm/arrow] (complete) -- node[farm/label,above] {Không} (supplement);
    \draw[farm/arrow] (supplement.north) |- (create.east);
    \draw[farm/arrow] (passed) -- node[farm/label,right] {Có} (trusted);
    \draw[farm/arrow] (passed) -- node[farm/label,above] {Không} (revise);
    \draw[farm/arrow] (revise) -- (reviseout);
    \draw[farm/arrow] (revisein.south) |- (create.east);
    \draw[farm/arrow] (trusted) -- node[farm/label,right] {Có} (publish);
    \draw[farm/arrow] (trusted) -- node[farm/label,above] {Không} (manual);
    \draw[farm/arrow] (accepted) -- node[farm/label,above] {Có} (publish);
    \draw[farm/arrow] (accepted.east) -- (8.6,-11.8) -- (8.6,-8.5) -| node[farm/label,near start,right] {Không} (revise.south);
    \draw[farm/arrow] (changed) -- node[farm/label,right] {Không} (finish);
    \draw[farm/arrow] (changed) -- node[farm/label,above] {Có} (pending);
    \draw[farm/arrow] (pending) -- (reviewout);
    \draw[farm/arrow] (reviewin) -- (automatic);
\end{tikzpicture}
\caption{Flowchart đăng và kiểm duyệt sản phẩm, sinh mã QR truy xuất.}
\label{fig:flow-san-pham}
\end{figure}
\clearpage

\clearpage
\subsubsection*{Quy trình bán lẻ (B2C)}
Quy trình bán lẻ được minh họa qua hai lưu đồ nối tiếp: Hình~\ref{fig:flow-ban-le-a} (đặt hàng, thanh toán và xác nhận đơn) và Hình~\ref{fig:flow-ban-le-b} (giao hàng, nhận hàng và hoàn tất đơn).
\begin{figure}[H]
\centering
\begin{adjustbox}{max totalsize={\textwidth}{\dimexpr\textheight-4.6cm\relax},center}
\begin{tikzpicture}[farm/chart]
    \node[farm/terminal] (start) at (0,0) {Bắt đầu bán lẻ};
    \node[farm/process] (ready) at (0,-1.7) {Nhà bán: sản phẩm đã duyệt (mục~\ref{sec:quytrinhnghiepvu}), hiển thị công khai với giá bán lẻ và sản lượng còn lại theo thời gian thực};
    \node[farm/process] (browse) at (0,-3.75) {Người mua: tìm kiếm, lọc, xem chi tiết và truy xuất nguồn gốc};
    \node[farm/process] (cart) at (0,-5.7) {Người mua: thêm vào giỏ theo sản lượng còn lại; chọn khung giờ giao và phương thức thanh toán};
    \node[farm/decision] (payment) at (0,-8) {Thanh toán thành công\\hoặc chọn COD?};
    \node[farm/processnarrow] (wait) at (7,-8) {Thanh toán trực tuyến thất bại: giữ đơn chờ tối đa 15 phút, cho phép thanh toán lại};
    \node[farm/decision] (retry) at (7,-11.4) {Thành công\\trong 15 phút?};
    \node[farm/decision] (stock) at (0,-11) {Còn đủ\\sản lượng?};
    \node[farm/exception] (cancel) at (7,-13.5) {Hệ thống: hủy đơn, thông báo lý do; giải phóng sản lượng đã giữ nếu có};
    \node[farm/process] (confirm) at (0,-13.5) {Hệ thống: xác nhận đơn, tạm trừ sản lượng khỏi tồn của nhà bán và báo nhà bán chuẩn bị hàng};
    \node[farm/terminal] (cancelled) at (7,-15.5) {Đơn bị hủy};
    \node[farm/connector] (connA) at (0,-15.5) {A};
    \node[farm/processnarrow] (connAnote) at (2.9,-15.5) {Luồng đi tiếp tại vòng tròn \textbf{A} của Hình~\ref{fig:flow-ban-le-b} (giao nhận hàng)};
    \node[farm/note] at (2.7,-17.3) {\textbf{Đặt đồng thời:} ưu tiên đơn thanh toán/xác nhận trước khi phân bổ sản lượng; đơn còn lại bị hủy và thông báo hết hàng.};
    \foreach \source/\target in {start/ready,ready/browse,browse/cart,cart/payment,wait/retry,confirm/connA,cancel/cancelled}
        \draw[farm/arrow] (\source) -- (\target);
    \draw[farm/arrow,dashed] (connA) -- (connAnote);
    \draw[farm/arrow] (payment) -- node[farm/label,left,pos=0.2] {Có} (stock);
    \draw[farm/arrow] (payment) -- node[farm/label,above] {Không} (wait);
    % Thanh toán lại thành công: nhập vào luồng chính ngay trên ô "Còn đủ sản lượng?".
    \draw[farm/arrow] (retry.west) -- node[farm/label,above] {Có} (3.7,-11.4) -- (3.7,-9.6) -- (0,-9.6);
    \draw[farm/arrow] (retry) -- node[farm/label,right] {Không} (cancel);
    \draw[farm/arrow] (stock) -- node[farm/label,right] {Có} (confirm);
    \draw[farm/arrow] (stock.east) -- node[farm/label,above] {Không} (3,-11) -- (3,-13.5) -- (cancel.west);
\end{tikzpicture}
\end{adjustbox}
\caption[Lưu đồ quy trình bán lẻ B2C (phần 1)]{Lưu đồ quy trình bán lẻ B2C (phần 1): đặt hàng, thanh toán và xác nhận đơn}
\label{fig:flow-ban-le-a}
\end{figure}

\begin{figure}[H]
\centering
\begin{adjustbox}{max totalsize={\textwidth}{\dimexpr\textheight-4.6cm\relax},center}
\begin{tikzpicture}[farm/chart]
    \node[farm/connector] (connA) at (0,0) {A};
    \node[farm/processnarrow] (connAnote) at (3.4,0) {Tiếp nối vòng tròn \textbf{A} của Hình~\ref{fig:flow-ban-le-a} (đơn đã xác nhận)};
    \node[farm/process] (ship) at (0,-2) {Nhà bán: đóng gói hàng dễ hư hỏng, bàn giao vận chuyển};
    \node[farm/decision] (refused) at (0,-4.1) {Người nhận\\từ chối COD?};
    \node[farm/exception] (return) at (6.5,-4.1) {Hoàn hàng; xử lý phí vận chuyển hai chiều theo chính sách};
    \node[farm/process] (receive) at (0,-6.4) {Người mua: xác nhận nhận hàng, đánh giá sản phẩm và nhà bán trong thời hạn quy định};
    \node[farm/process] (sellerdone) at (0,-8.4) {Nhà bán: đơn hoàn tất --- hệ thống cập nhật doanh thu và lịch sử bán lẻ};
    \node[farm/terminal] (finish) at (0,-10.3) {Kết thúc};
    \node[farm/note] at (2.7,-12) {\textbf{Hàng hư hỏng:} người nhận khiếu nại và được xử lý theo flowchart vận chuyển và khiếu nại tại Hình~\ref{fig:flow-van-chuyen}.};
    \foreach \source/\target in {connA/ship,ship/refused,receive/sellerdone,sellerdone/finish}
        \draw[farm/arrow] (\source) -- (\target);
    \draw[farm/arrow,dashed] (connAnote) -- (connA);
    \draw[farm/arrow] (refused) -- node[farm/label,right] {Không} (receive);
    \draw[farm/arrow] (refused) -- node[farm/label,above] {Có} (return);
    \draw[farm/arrow] (return.south) |- (finish.east);
\end{tikzpicture}
\end{adjustbox}
\caption[Lưu đồ quy trình bán lẻ B2C (phần 2)]{Lưu đồ quy trình bán lẻ B2C (phần 2): giao hàng, nhận hàng và hoàn tất đơn}
\label{fig:flow-ban-le-b}
\end{figure}


\clearpage
\subsubsection*{Đánh giá sản phẩm và nhà bán}
Hình~\ref{fig:flow-danh-gia} minh họa quy trình đánh giá, từ kiểm duyệt đánh giá tới cập nhật điểm uy tín và quyền lợi của nhà bán.
\clearpage
\begin{figure}[p]
\centering
\resizebox{!}{\dimexpr\textheight-2cm\relax}{%
\begin{tikzpicture}[farm/chart]
    \node[farm/terminal] (start) at (0,0) {Đơn hàng bán lẻ hoàn tất (verified purchase)};
    \node[farm/decision] (window) at (0,-2) {Trong thời hạn\\đánh giá quy định?};
    \node[farm/exception] (locked) at (6,-2) {Hệ thống: khóa chức năng đánh giá cho đơn này};
    \node[farm/process] (submit) at (0,-4.2) {Người mua: gửi đánh giá (sao 1--5, bình luận, ảnh minh chứng)};
    \node[farm/process] (autocheck) at (0,-6) {Hệ thống: kiểm duyệt tự động sơ bộ + mô-đun AI phát hiện bất thường};
    \node[farm/decision] (suspect) at (0,-8.2) {Nghi ngờ\\giả/spam?};
    \node[farm/process] (manual) at (6,-8.2) {Quản trị viên: xem xét thủ công};
    \node[farm/decision] (approved) at (6,-10.4) {Được duyệt?};
    \node[farm/exception] (reject) at (9.5,-10.4) {Ẩn đánh giá, thông báo lý do cho người mua};
    \node[farm/process] (publish) at (0,-10.4) {Hệ thống: công khai đánh giá};
    \node[farm/process] (score) at (0,-12.2) {Hệ thống: cập nhật điểm uy tín tổng hợp của nhà bán};
    \node[farm/process] (reply) at (0,-14) {Nhà bán: có thể phản hồi công khai (một lần) dưới đánh giá};
    \node[farm/decision] (threshold) at (0,-16.2) {Điểm uy tín\\đủ ngưỡng?};
    \node[farm/side] (perks) at (6,-16.2) {Nhà bán: đủ điều kiện hậu kiểm nhanh và hiển thị ưu tiên};
    \node[farm/terminal] (finish) at (0,-18.4) {Kết thúc};
    \node[farm/note] at (2.7,-20.2) {\textbf{Điểm uy tín:} trung bình có trọng số theo thời gian và số lượng đơn, không tính đánh giá bị ẩn.\\[2pt]\textbf{Liên hệ:} mô-đun AI phát hiện bất thường tại mục~\ref{sec:thietkeAI}; điều kiện hậu kiểm nhanh tại mục~\ref{sec:quytrinhnghiepvu} (Đăng và kiểm duyệt sản phẩm).};
    \foreach \source/\target in {start/window,submit/autocheck,autocheck/suspect,score/reply,reply/threshold}
        \draw[farm/arrow] (\source) -- (\target);
    \draw[farm/arrow] (window) -- node[farm/label,right] {Có} (submit);
    \draw[farm/arrow] (window) -- node[farm/label,above] {Không} (locked);
    \draw[farm/arrow] (locked.south) |- (finish.east);
    \draw[farm/arrow] (suspect) -- node[farm/label,right] {Không} (publish);
    \draw[farm/arrow] (suspect) -- node[farm/label,above] {Có} (manual);
    \draw[farm/arrow] (manual) -- (approved);
    \draw[farm/arrow] (approved) -- node[farm/label,above] {Không} (reject);
    \draw[farm/arrow] (approved.west) -- node[farm/label,near start,above] {Có} (publish.east);
    \draw[farm/arrow] (reject.south) |- (finish.east);
    \draw[farm/arrow] (publish) -- (score);
    \draw[farm/arrow] (threshold) -- node[farm/label,right] {Có} (perks);
    \draw[farm/arrow] (threshold) -- node[farm/label,above] {Không} (finish);
    \draw[farm/arrow] (perks.south) |- (finish.east);
\end{tikzpicture}%
}
\caption{Flowchart đánh giá sản phẩm và nhà bán: kiểm duyệt, cập nhật điểm uy tín và tác động tới quyền lợi nhà bán.}
\label{fig:flow-danh-gia}
\end{figure}
\clearpage


\clearpage
\subsubsection*{Quy trình bán sỉ (B2B)}
Quy trình bán sỉ được minh họa qua hai lưu đồ nối tiếp: Hình~\ref{fig:flow-ban-si-a} (từ công bố bảng giá tới ký hợp đồng) và Hình~\ref{fig:flow-ban-si-b} (ký quỹ, giao hàng theo đợt và công nợ).
\clearpage
\begin{figure}[p]
\centering
\resizebox{!}{\dimexpr\textheight-2cm\relax}{%
\begin{tikzpicture}[farm/chart]
    \node[farm/terminal] (start) at (0,0) {Bắt đầu bán sỉ};
    \node[farm/process] (pricelist) at (0,-1.4) {Nhà bán: công bố bảng giá bán sỉ theo bậc số lượng, MOQ và lịch giao dự kiến cho từng lô hàng};
    \node[farm/process] (rfq) at (0,-2.9) {Doanh nghiệp: xem bảng giá công khai, gửi RFQ tham chiếu bảng giá hoặc yêu cầu tùy chỉnh (số lượng, chất lượng, lịch giao)};
    \node[farm/decision] (moq) at (0,-4.8) {Đạt MOQ\\của nhà bán?};
    \node[farm/side] (alternative) at (6,-4.8) {Hệ thống: gợi ý gộp đơn cùng khu vực hoặc chuyển kênh bán lẻ};
    \node[farm/terminal] (retail) at (6,-6.8) {Chuyển sang bán lẻ (B2C)};
    \node[farm/process] (quote) at (0,-6.8) {Nhà bán: xác nhận/điều chỉnh báo giá theo bậc, thời gian đáp ứng và điều kiện thanh toán};
    \node[farm/process] (negotiate) at (0,-8.3) {Hai bên: thương lượng qua hệ thống nhắn tin nội bộ};
    \node[farm/decision] (agreed) at (0,-10.1) {Thống nhất\\điều khoản?};
    \node[farm/process] (contract) at (0,-12.1) {Hệ thống: sinh hợp đồng, gồm phạt vi phạm và công nợ; hai bên ký điện tử / mã xác nhận};
    \node[farm/process] (delivery) at (0,-13.9) {Nhà bán giao theo từng đợt; hệ thống lập hóa đơn và cập nhật sổ công nợ};
    \node[farm/process] (settle) at (0,-15.6) {Doanh nghiệp: xác nhận số lượng, chất lượng từng đợt; thanh toán theo kỳ hạn};
    \node[farm/decision] (complete) at (0,-17.5) {Còn đợt giao\\theo hợp đồng?};
    \node[farm/side] (reconcile) at (6,-17.5) {Hai bên: đối soát, hoàn tất thanh toán công nợ theo kỳ hạn};
    \node[farm/process] (sellerdone) at (0,-19.5) {Nhà bán: hợp đồng hoàn tất — hệ thống cập nhật doanh thu và lịch sử bán sỉ của nhà bán};
    \node[farm/terminal] (finish) at (0,-21.1) {Hoàn tất hợp đồng};
    \node[farm/exception,text width=3.8cm] (exceptions) at (6,-12.1) {\textbf{Ngoại lệ hợp đồng}\\[3pt]Giao thiếu / trễ: áp dụng phạt và ghi lịch sử uy tín.\\[4pt]Quá hạn công nợ: nhắc tự động; có thể khóa RFQ mới đến khi trả đủ.\\[4pt]Hủy giữa kỳ: xử lý theo điều khoản hủy đã ký.};
    \foreach \source/\target in {start/pricelist,pricelist/rfq,rfq/moq,quote/negotiate,negotiate/agreed,contract/delivery,delivery/settle,settle/complete,sellerdone/finish}
        \draw[farm/arrow] (\source) -- (\target);
    \draw[farm/arrow] (moq) -- node[farm/label,right] {Có} (quote);
    \draw[farm/arrow] (moq) -- node[farm/label,above] {Không} (alternative);
    \draw[farm/arrow] (alternative.north) |- node[farm/label,near start,right] {Gộp đơn, gửi lại} (rfq.east);
    \draw[farm/arrow] (alternative) -- node[farm/label,right] {Chuyển kênh} (retail);
    \draw[farm/arrow] (agreed) -- node[farm/label,right] {Có} (contract);
    \draw[farm/arrow] (agreed.west) -- (-3.5,-10.1) |- node[farm/label,near start,left] {Chưa} (quote.west);
    \draw[farm/arrow] (complete) -- node[farm/label,above] {Không} (reconcile);
    \draw[farm/arrow] (reconcile.south) |- (sellerdone.east);
    \draw[farm/arrow] (complete.west) -- (-3.5,-17.5) |- node[farm/label,near start,left] {Có} (delivery.west);
    \draw[farm/arrow,dashed] (contract) -- (exceptions);
\end{tikzpicture}%
}
\caption{Flowchart bán sỉ (B2B): công bố bảng giá, báo giá/RFQ, hợp đồng, giao hàng và công nợ.}
\label{fig:flow-ban-si}
\end{figure}
\clearpage

\begin{figure}[H]
\centering
\begin{adjustbox}{max totalsize={\textwidth}{\dimexpr\textheight-4.6cm\relax},center}
\begin{tikzpicture}[farm/chart]
    \node[farm/connector] (connA) at (0,0) {A};
    \node[farm/processnarrow] (connAnote) at (3.9,0) {Tiếp nối vòng tròn \textbf{A} của Hình~\ref{fig:flow-ban-si-a} (hợp đồng vừa được ký)};
    \draw[farm/arrow,dashed] (connAnote) -- (connA);
    \node[farm/decision] (escrowq) at (0,-2) {Giá trị vượt ngưỡng\\ký quỹ (FR7.4)?};
    \node[farm/processnarrow] (nonescrow) at (6,-2) {Không ký quỹ: thanh toán theo kỳ hạn công nợ đã thỏa thuận};
    \node[farm/process] (escrow) at (0,-4.8) {Bên mua nộp tiền đợt giao vào tài khoản ký quỹ tại cổng thanh toán trung gian; hệ thống ghi nhận trạng thái ``đã ký quỹ''};
    \node[farm/process] (delivery) at (0,-7.2) {Nhà bán giao theo từng đợt; hệ thống lập hóa đơn và cập nhật sổ công nợ};
    \node[farm/process] (settle) at (0,-9.2) {Doanh nghiệp: xác nhận số lượng và chất lượng của đợt giao};
    \node[farm/process] (release) at (0,-11.2) {Hệ thống: ghi nhận thanh toán; nếu có ký quỹ thì ra lệnh cổng trung gian giải ngân khoản đã giữ cho nhà bán};
    \node[farm/decision] (complete) at (0,-13.6) {Còn đợt giao\\theo hợp đồng?};
    \node[farm/processnarrow] (reconcile) at (6,-13.6) {Hai bên: đối soát, hoàn tất thanh toán công nợ theo kỳ hạn};
    \node[farm/process] (sellerdone) at (0,-15.8) {Nhà bán: hợp đồng hoàn tất --- hệ thống cập nhật doanh thu và lịch sử bán sỉ};
    \node[farm/terminal] (finish) at (0,-17.6) {Hoàn tất hợp đồng};
    \node[farm/exception,text width=3.8cm] (exceptions) at (6,-8.1) {\textbf{Ngoại lệ hợp đồng}\\[3pt]Giao thiếu / trễ: áp dụng phạt và ghi lịch sử uy tín.\\[4pt]Quá hạn công nợ: nhắc tự động; có thể khóa RFQ mới đến khi trả đủ.\\[4pt]Hủy giữa kỳ: xử lý theo điều khoản hủy đã ký.\\[4pt]Tranh chấp đợt đã ký quỹ: giữ nguyên khoản ký quỹ tại cổng trung gian đến khi xử lý xong khiếu nại.};
    \foreach \source/\target in {connA/escrowq,escrow/delivery,delivery/settle,settle/release,release/complete,sellerdone/finish}
        \draw[farm/arrow] (\source) -- (\target);
    \draw[farm/arrow] (escrowq) -- node[farm/label,right] {Có} (escrow);
    \draw[farm/arrow] (escrowq) -- node[farm/label,above] {Không} (nonescrow);
    % Đi vòng qua rãnh trống x=3.3 để không cắt ngang hộp ngoại lệ đặt tại x=6.
    \draw[farm/arrow] (nonescrow.south) -- (6,-3.05) -- (3.3,-3.05) -- (3.3,-7.2) -- (delivery.east);
    \draw[farm/arrow] (complete) -- node[farm/label,above] {Không} (reconcile);
    \draw[farm/arrow] (reconcile.south) |- (sellerdone.east);
    \draw[farm/arrow] (complete.west) -- (-3.5,-13.6) |- node[farm/label,near start,left] {Có} (delivery.west);
    \draw[farm/arrow,dashed] (delivery) -- (exceptions);
\end{tikzpicture}
\end{adjustbox}
\caption[Lưu đồ quy trình bán sỉ B2B (phần 2)]{Lưu đồ quy trình bán sỉ B2B (phần 2): ký quỹ, giao hàng theo đợt, giải ngân và công nợ}
\label{fig:flow-ban-si-b}
\end{figure}


\clearpage
\subsubsection*{Quản lý tồn kho, vận chuyển hàng dễ hư hỏng và xử lý khiếu nại}
Hình~\ref{fig:flow-van-chuyen} minh họa quy trình vận chuyển hàng dễ hư hỏng và xử lý khiếu nại.
\clearpage
\begin{figure}[p]
\centering
\begin{adjustbox}{max totalsize={\textwidth}{\dimexpr\textheight-2.4cm\relax},center}
\begin{tikzpicture}[farm/chart]
    \node[farm/terminal] (start) at (0,0) {Đơn hàng dự kiến};
    \node[farm/process] (method) at (0,-1.35) {Hệ thống: ưu tiên giao nhanh / bảo quản lạnh theo danh mục dễ hư hỏng};
    \node[farm/decision] (unsafe) at (0,-3.2) {Thời gian giao vượt ngưỡng an toàn?};
    \node[farm/processnarrow] (warn) at (6,-3.2) {Hệ thống: cảnh báo trước khi khách xác nhận đặt hàng};
    \node[farm/process] (confirm) at (0,-5.6) {Khách xác nhận đặt hàng;\\nhà bán đóng gói theo nhiệt độ, độ ẩm và vật liệu phù hợp};
    \node[farm/process] (transport) at (0,-7.3) {Đơn vị vận chuyển: giao hàng, cập nhật trạng thái theo thời gian thực};
    \node[farm/decision] (damage) at (0,-9.4) {Người nhận kiểm tra: hàng hư hỏng?};
    \node[farm/terminal] (received) at (6,-9.4) {Nhận hàng, hoàn tất};
    \node[farm/process] (claim) at (0,-11.4) {Người nhận: khiếu nại trong thời hạn, kèm ảnh / video minh chứng};
    \node[farm/decision] (suspicious) at (0,-13.4) {Người mua có\\dấu hiệu bất thường?};
    \node[farm/processnarrow] (manual) at (6,-13.4) {Quản trị: xem xét thủ công trước khi quyết định hoàn tiền};
    \node[farm/process] (classify) at (0,-15.8) {Hệ thống: phân loại nguyên nhân từ nhà bán, vận chuyển hoặc bảo quản của người nhận};
    \node[farm/process] (resolve) at (0,-17.7) {Bộ phận phụ trách: hoàn tiền toàn phần / một phần, đổi lô mới hoặc từ chối nếu thiếu bằng chứng};
    \node[farm/terminal] (finish) at (0,-19.4) {Thông báo kết quả xử lý};
    \node[farm/exception] (repeated) at (6,-17.6) {Khiếu nại lặp lại từ cùng nhà bán trong thời gian ngắn:\\gắn cờ và tạm hạ ưu tiên hiển thị đến khi rà soát};
    \foreach \source/\target in {start/method,method/unsafe,confirm/transport,transport/damage,claim/suspicious,classify/resolve,resolve/finish}
        \draw[farm/arrow] (\source) -- (\target);
    \draw[farm/arrow] (unsafe) -- node[farm/label,right] {Không} (confirm);
    \draw[farm/arrow] (unsafe) -- node[farm/label,above] {Có} (warn);
    \draw[farm/arrow] (warn.south) |- (confirm.east);
    \draw[farm/arrow] (damage) -- node[farm/label,right] {Có} (claim);
    \draw[farm/arrow] (damage) -- node[farm/label,above] {Không} (received);
    \draw[farm/arrow] (suspicious) -- node[farm/label,right] {Không} (classify);
    \draw[farm/arrow] (suspicious) -- node[farm/label,above] {Có} (manual);
    \draw[farm/arrow] (manual.south) |- (classify.east);
    \draw[farm/arrow,dashed] (resolve) -- (repeated);
\end{tikzpicture}
\end{adjustbox}
\caption{Flowchart vận chuyển hàng dễ hư hỏng và xử lý khiếu nại.}
\label{fig:flow-van-chuyen}
\end{figure}
\clearpage

\clearpage
\subsubsection*{Quản trị hệ thống}
Hình~\ref{fig:flow-quan-tri} minh họa quy trình giám sát, xử lý vi phạm phân cấp và báo cáo của quản trị viên.
\begin{figure}[H]
\centering
\begin{adjustbox}{max totalsize={\textwidth}{\dimexpr\textheight-4.6cm\relax},center}
\begin{tikzpicture}[farm/chart]
    \node[farm/terminal] (start) at (0,0) {Bắt đầu đợt giám sát};
    \node[farm/process] (review) at (0,-1.5) {Quản trị: kiểm duyệt hồ sơ mới, sản phẩm mới / thay đổi trọng yếu và chứng nhận sắp hết hạn};
    \node[farm/process] (monitor) at (0,-3.2) {Quản trị: theo dõi cảnh báo tỷ lệ khiếu nại bất thường và giao dịch có nguy cơ gian lận};
    \node[farm/decision] (violation) at (0,-5.2) {Phát hiện\\vi phạm?};
    \node[farm/decision] (severe) at (0,-7.9) {Vi phạm\\nghiêm trọng?};
    \node[farm/exception] (lock) at (6,-7.9) {Giả mạo chứng nhận / gian lận:\\khóa ngay tài khoản, không qua nhắc nhở trung gian};
    \node[farm/process] (sanction) at (0,-10.5) {Quản trị: xử lý theo mức độ và số lần vi phạm: nhắc nhở, cảnh cáo công khai, tạm khóa hoặc khóa vĩnh viễn};
    \node[farm/exception] (orders) at (6,-10.5) {Rà soát toàn bộ đơn đang xử lý của tài khoản, bảo vệ khách đã đặt trước};
    \node[farm/process] (audit) at (0,-12.7) {Hệ thống / quản trị: lưu vết đầy đủ lý do và quyết định xử lý};
    \node[farm/process] (report) at (0,-14.5) {Quản trị: truy xuất báo cáo định kỳ, hỗ trợ quyết định vận hành và báo cáo các bên liên quan};
    \node[farm/terminal] (finish) at (0,-16.2) {Kết thúc đợt giám sát};
    \node[farm/note] at (2.7,-18) {\textbf{Nguyên tắc phân cấp:} các hình thức xử lý được lựa chọn theo mức độ và số lần vi phạm, không bắt buộc đi qua mọi mức. Quy trình giám sát được thực hiện định kỳ; trường hợp nghiêm trọng được xử lý ngay khi phát hiện.};
    \foreach \source/\target in {start/review,review/monitor,monitor/violation,lock/orders,sanction/audit,audit/report,report/finish}
        \draw[farm/arrow] (\source) -- (\target);
    \draw[farm/arrow] (violation) -- node[farm/label,right] {Có} (severe);
    \draw[farm/arrow] (violation.west) -- (-3.6,-5.2) |- node[farm/label,near start,left] {Không} (report.west);
    \draw[farm/arrow] (severe) -- node[farm/label,right] {Không} (sanction);
    \draw[farm/arrow] (severe) -- node[farm/label,above] {Có} (lock);
    \draw[farm/arrow] (orders.south) |- (audit.east);
\end{tikzpicture}
\end{adjustbox}
\caption{Lưu đồ quy trình quản trị hệ thống, xử lý vi phạm và báo cáo}
\label{fig:flow-quan-tri}
\end{figure}

\clearpage
\subsection{Mô hình hóa tương tác}
\label{sec:sddchitiet}

% TODO: vẽ use case diagram tổng quát và theo từng tác nhân (nhà bán, Nhóm B, Nhóm C, các vai trò quản trị viên tại mục 4.2.1); kèm bảng mô tả cho từng use case chính (tên, tác nhân, tiền điều kiện, luồng chính, luồng thay thế/ngoại lệ, hậu điều kiện), bám theo mã FR tại mục 4.3.1 và các quy trình tại mục 4.1.2.

\section{Giải pháp công nghệ}
\label{sec:giaiphapcongnghe}

% TODO: với mỗi hạng mục, liệt kê các phương án ứng viên, tiêu chí so sánh (hiệu năng, hệ sinh thái, độ quen thuộc của nhóm, chi phí, khả năng đáp ứng NFR tại Bảng 4.x) và lý do chọn phương án cuối cùng.

\subsection{Front-end}
% TODO

\subsection{Back-end}
% TODO: gồm framework, hệ quản trị CSDL, lưu trữ đối tượng, cổng thanh toán và đơn vị vận chuyển tích hợp.

\subsection{Thuật toán, công nghệ và mô hình AI}
% TODO: lựa chọn mô hình/thư viện cho 4 mô-đun AI tại mục 5.3.6, đối chiếu với ngưỡng chất lượng tại Bảng~\ref{tab:nguongAI}.

\section{Thiết kế}
\label{sec:thietke}

\subsection{Thiết kế kiến trúc hệ thống}
\label{sec:kientruc}

% TODO: nhóm bổ sung sơ đồ kiến trúc cụ thể (\includegraphics) và mô tả rõ công nghệ dự kiến/đã chọn cho từng lớp (framework, ngôn ngữ, hệ quản trị CSDL...).
Áp dụng mô hình kiến trúc phân lớp đã trình bày tại Chương~\ref{chap:coso} vào Farmery: sơ đồ kiến trúc cụ thể (thành phần, công nghệ từng lớp, luồng gọi giữa các lớp) sẽ được bổ sung ở đây sau khi nhóm chốt giải pháp công nghệ tại mục~\ref{sec:giaiphapcongnghe}.


\subsection{Thiết kế sitemap}
% TODO: sơ đồ các trang/màn hình theo từng nhóm người dùng.

\subsection{Thiết kế lược đồ tuần tự}
% TODO: sequence diagram cho các luồng chính: đặt hàng B2C, RFQ -- hợp đồng -- ký quỹ B2B, đăng và kiểm duyệt sản phẩm, gửi và kiểm duyệt đánh giá.

\subsection{Thiết kế lược đồ lớp}
% TODO: class diagram cho các lớp miền chính, bám theo các nhóm thực thể tại mục 5.3.5.

\subsection{Thiết kế cơ sở dữ liệu}
\label{sec:csdl}

\subsubsection*{Đặc tả cơ sở dữ liệu}
\label{subsec:dacta-csdl}

Cơ sở dữ liệu được tổ chức theo 8 nhóm thực thể tương ứng với các nhóm chức
năng FR1--FR12 tại mục~\ref{sec:yeucauchucnang}, nhằm đảm bảo khả năng tách
module theo domain (NFR5.1). Quy ước: \textbf{PK} = khóa chính,
\textit{FK} = khóa ngoại. Bốn nhóm thực thể cốt lõi (Nhóm 1--4) được đặc tả
dạng bảng trường--kiểu--mô tả để tiện tra cứu khi hiện thực; các nhóm còn lại
(Nhóm 5--8) có số trường ít hơn nên giữ dạng danh sách.
 
\paragraph*{Nhóm 1 --- Người dùng \& phân quyền (ứng với FR1)}

\noindent Bảng~\ref{tab:db-nhom1} trình bày đặc tả các thực thể thuộc nhóm này.
 
\begingroup
\small
\renewcommand{\arraystretch}{1.2}
\setlength{\tabcolsep}{5pt}
\begin{longtable}{L{4.9cm} L{3.8cm} L{6.0cm}}
\caption{Đặc tả thực thể nhóm 1: người dùng và phân quyền}
\label{tab:db-nhom1} \\
\toprule
\rowcolor{tblheadbg}
\textbf{Trường} & \textbf{Kiểu / Ràng buộc} & \textbf{Mô tả \& FR liên quan} \\
\midrule
\endfirsthead
\toprule
\rowcolor{tblheadbg}
\textbf{Trường} & \textbf{Kiểu / Ràng buộc} & \textbf{Mô tả \& FR liên quan} \\
\midrule
\endhead
\bottomrule
\endfoot
\multicolumn{3}{@{}l}{\cellcolor{tblaltbg}\textbf{User} (Người dùng)} \\
\texttt{id} & PK & Định danh người dùng. \\
\texttt{full\_name} & & Họ tên hiển thị. \\
\texttt{email} & unique, nullable & Kênh đăng nhập và nhận thông báo. \\
\texttt{phone} & unique, nullable & Kênh đăng nhập và nhận OTP (FR1.2). \\
\texttt{password\_hash} & nullable & Băm bằng Bcrypt/Argon2 (NFR3.3); rỗng khi tài khoản chỉ đăng nhập qua mạng xã hội. \\
\texttt{oauth\_provider} & nullable & Nhà cung cấp OAuth (Google, Facebook...) --- FR1.1. \\
\texttt{oauth\_id} & nullable & Định danh tài khoản mạng xã hội tương ứng --- FR1.1. \\
\texttt{role} & enum & \emph{seller, buyer\_retail, buyer\_wholesale, admin, admin\_moderator, admin\_compliance, admin\_support, admin\_ops, admin\_inventory} --- trùng đúng tên vai trò RBAC tại mục~\ref{sec:doituongnguoidung}. \\
\texttt{status} & enum & Hoạt động / khóa. \\
\texttt{created\_at} & datetime & Thời điểm tạo tài khoản. \\
\midrule
\multicolumn{3}{@{}l}{\cellcolor{tblaltbg}\textbf{Address} (Địa chỉ) --- sổ địa chỉ giao/nhận hàng, FR1.6} \\
\texttt{id} & PK & \\
\texttt{user\_id} & FK $\to$ User.id & Chủ sở hữu địa chỉ. \\
\texttt{content} & & Địa chỉ đầy đủ. \\
\texttt{label} & & Nhãn gợi nhớ (nhà riêng, công ty...). \\
\texttt{is\_default} & boolean & Địa chỉ mặc định khi đặt hàng. \\
\midrule
\multicolumn{3}{@{}l}{\cellcolor{tblaltbg}\textbf{Business} (Doanh nghiệp) --- hồ sơ doanh nghiệp thu mua, FR1.3} \\
\texttt{id} & PK & \\
\texttt{owner\_user\_id} & FK $\to$ User.id & Người đứng tên tài khoản doanh nghiệp. \\
\texttt{company\_name} & & \\
\texttt{tax\_code} & unique & Mã số thuế dùng để xác minh. \\
\texttt{business\_license\_file} & & Tệp giấy phép kinh doanh đính kèm. \\
\texttt{verification\_status} & enum & Chờ duyệt / đã duyệt / từ chối. \\
\midrule
\multicolumn{3}{@{}l}{\cellcolor{tblaltbg}\textbf{BusinessMember} (Thành viên doanh nghiệp) --- RBAC nội bộ, FR1.4} \\
\texttt{id} & PK & \\
\texttt{business\_id} & FK $\to$ Business.id & \\
\texttt{user\_id} & FK $\to$ User.id & \\
\texttt{internal\_role} & enum & \emph{owner, buyer, accountant, viewer}. \\
\end{longtable}
\endgroup
Quan hệ: một \emph{User} có nhiều \emph{Address} (1--n); một
\emph{Business} có nhiều \emph{BusinessMember}, mỗi thành viên gắn
với một \emph{User} (n--n giữa User và Business, giải quyết qua
bảng trung gian BusinessMember).
 
\paragraph*{Nhóm 2 --- Gian hàng, chứng nhận, sản phẩm \& truy xuất nguồn gốc
(ứng với FR2--FR4)}

\noindent Bảng~\ref{tab:db-nhom2} trình bày đặc tả các thực thể thuộc nhóm này.
 
\begingroup
\small
\renewcommand{\arraystretch}{1.2}
\setlength{\tabcolsep}{5pt}
\begin{longtable}{L{4.9cm} L{3.8cm} L{6.0cm}}
\caption[Đặc tả thực thể nhóm 2: gian hàng, sản phẩm và truy xuất nguồn gốc]{Đặc tả thực thể nhóm 2: gian hàng, chứng nhận, sản phẩm và truy xuất nguồn gốc}
\label{tab:db-nhom2} \\
\toprule
\rowcolor{tblheadbg}
\textbf{Trường} & \textbf{Kiểu / Ràng buộc} & \textbf{Mô tả \& FR liên quan} \\
\midrule
\endfirsthead
\toprule
\rowcolor{tblheadbg}
\textbf{Trường} & \textbf{Kiểu / Ràng buộc} & \textbf{Mô tả \& FR liên quan} \\
\midrule
\endhead
\bottomrule
\endfoot
\multicolumn{3}{@{}l}{\cellcolor{tblaltbg}\textbf{Shop} (Gian hàng) --- FR2.1} \\
\texttt{id} & PK & \\
\texttt{user\_id} & FK $\to$ User.id & Ràng buộc \texttt{role = seller}. \\
\texttt{shop\_name} / \texttt{description} & & Thông tin hiển thị của gian hàng. \\
\texttt{growing\_region} & & Khu vực canh tác công bố. \\
\texttt{ekyc\_status} & enum & Chờ duyệt / đã duyệt / từ chối --- FR1.5. \\
\texttt{reputation\_score} & 0--100 & Giá trị dẫn xuất, tính lại sau mỗi đánh giá hợp lệ theo công thức tại mục~\ref{subsec:danhgianghiepvu}. \\
\texttt{reputation\_updated\_at} & datetime & Lần cập nhật điểm uy tín gần nhất. \\
\midrule
\multicolumn{3}{@{}l}{\cellcolor{tblaltbg}\textbf{Certification} (Chứng nhận) --- FR3.1--FR3.3} \\
\texttt{id} & PK & \\
\texttt{shop\_id} & FK $\to$ Shop.id & \\
\texttt{type} & enum & \emph{VietGAP, GlobalGAP, organic, OCOP, HACCP} --- đúng danh sách tại FR3.1. \\
\texttt{certificate\_number} & & Mã số do cơ quan cấp phát hành. \\
\texttt{issuing\_authority} & & Tổ chức cấp chứng nhận. \\
\texttt{evidence\_file} & & Tệp bản scan chứng nhận. \\
\texttt{issued\_date} / \texttt{expiry\_date} & date & Theo dõi hiệu lực, nhắc gia hạn (FR3.3). \\
\texttt{approval\_status} & enum & Chờ / đã duyệt / từ chối (FR3.2). \\
\texttt{reviewer\_id} & FK $\to$ User.id, nullable & Quản trị viên kiểm duyệt. \\
\midrule
\multicolumn{3}{@{}l}{\cellcolor{tblaltbg}\textbf{Category} (Danh mục) --- phân loại đa cấp, FR4.4} \\
\texttt{id} & PK & \\
\texttt{name} & & Tên danh mục. \\
\texttt{parent\_category\_id} & FK $\to$ Category.id, nullable & Tự tham chiếu để tạo cây danh mục. \\
\midrule
\multicolumn{3}{@{}l}{\cellcolor{tblaltbg}\textbf{Product} (Sản phẩm) --- FR4.1} \\
\texttt{id} & PK & \\
\texttt{shop\_id} & FK $\to$ Shop.id & Gian hàng sở hữu sản phẩm. \\
\texttt{category\_id} & FK $\to$ Category.id & \\
\texttt{name} / \texttt{description} & & Thông tin hiển thị. \\
\texttt{unit} & & kg, tấn, thùng... \\
\texttt{retail\_price} & & Giá bán lẻ B2C. \\
\texttt{status} & enum & Đang bán / ẩn / chờ duyệt. \\
\midrule
\multicolumn{3}{@{}l}{\cellcolor{tblaltbg}\textbf{WholesalePriceTier} (Bảng giá sỉ) --- giá bậc theo MOQ, FR4.1 và FR6.4} \\
\texttt{id} & PK & \\
\texttt{product\_id} & FK $\to$ Product.id & \\
\texttt{min\_quantity} & & Ngưỡng số lượng bắt đầu áp dụng bậc giá. \\
\texttt{unit\_price} & & Đơn giá tương ứng bậc. \\
\midrule
\multicolumn{3}{@{}l}{\cellcolor{tblaltbg}\textbf{Batch} (Lô hàng) --- đơn vị truy xuất nguồn gốc, FR2.3 và FR4.2} \\
\texttt{id} & PK & \\
\texttt{product\_id} & FK $\to$ Product.id & \\
\texttt{batch\_code} & unique & Mã lô hàng. \\
\texttt{harvest\_date} / \texttt{expiry\_date} & date & Cơ sở áp dụng nguyên tắc FEFO (mục~\ref{subsec:khobai}). \\
\texttt{stock\_quantity} & & Tồn kho hiện tại của lô. \\
\texttt{qr\_code} & unique & Mã QR truy xuất gắn với lô. \\
\midrule
\multicolumn{3}{@{}l}{\cellcolor{tblaltbg}\textbf{FarmingLog} (Nhật ký canh tác) --- FR4.2} \\
\texttt{id} & PK & \\
\texttt{batch\_id} & FK $\to$ Batch.id & \\
\texttt{activity\_date} & date & Ngày thực hiện hoạt động canh tác. \\
\texttt{activity} / \texttt{description} & & Gieo trồng, chăm sóc, thu hoạch... \\
\texttt{images} & & Ảnh minh chứng. \\
\texttt{is\_published} & boolean & Đã hiển thị trên trang truy xuất hay chưa. \\
\end{longtable}
\endgroup
Quan hệ: Shop 1--n Certification; Shop 1--n Product; Product 1--n
WholesalePriceTier; Product 1--n Batch; Batch 1--n FarmingLog.
\textbf{Ràng buộc toàn vẹn (NFR9.1):} khi \texttt{Certification.approval\_status
= approved} hoặc \texttt{FarmingLog.is\_published = true}, các trường nội
dung không được phép \texttt{UPDATE} trực tiếp --- mọi chỉnh sửa sau công khai
phải tạo bản ghi mới và giữ liên kết lịch sử (append-only), liên hệ
mục~\ref{subsec:gs1blockchain}.
 
\paragraph*{Nhóm 3 --- Giao dịch bán lẻ B2C (ứng với FR5)}

\noindent Bảng~\ref{tab:db-nhom3} trình bày đặc tả các thực thể thuộc nhóm này.
 
\begingroup
\small
\renewcommand{\arraystretch}{1.2}
\setlength{\tabcolsep}{5pt}
\begin{longtable}{L{4.9cm} L{3.8cm} L{6.0cm}}
\caption{Đặc tả thực thể nhóm 3: giao dịch bán lẻ B2C}
\label{tab:db-nhom3} \\
\toprule
\rowcolor{tblheadbg}
\textbf{Trường} & \textbf{Kiểu / Ràng buộc} & \textbf{Mô tả \& FR liên quan} \\
\midrule
\endfirsthead
\toprule
\rowcolor{tblheadbg}
\textbf{Trường} & \textbf{Kiểu / Ràng buộc} & \textbf{Mô tả \& FR liên quan} \\
\midrule
\endhead
\bottomrule
\endfoot
\multicolumn{3}{@{}l}{\cellcolor{tblaltbg}\textbf{Cart} (Giỏ hàng) \textbf{/ CartItem} (Chi tiết giỏ hàng) --- FR5.1} \\
\texttt{Cart.id} & PK & Mỗi người mua một giỏ hàng. \\
\texttt{Cart.user\_id} & FK $\to$ User.id & \\
\texttt{CartItem.id} & PK & \\
\texttt{cart\_id} & FK $\to$ Cart.id & \\
\texttt{product\_id} & FK $\to$ Product.id & Sản phẩm trong giỏ, có thể thuộc nhiều gian hàng. \\
\texttt{quantity} & & \\
\midrule
\multicolumn{3}{@{}l}{\cellcolor{tblaltbg}\textbf{Voucher} (Mã giảm giá) --- FR5.3} \\
\texttt{id} & PK & \\
\texttt{code} & unique & Mã người mua nhập khi thanh toán. \\
\texttt{scope} & enum & \emph{platform} hoặc \emph{shop}. \\
\texttt{shop\_id} & FK $\to$ Shop.id, nullable & Chỉ dùng khi phạm vi là gian hàng. \\
\texttt{discount\_percent} / \texttt{max\_discount} & & Mức giảm và trần giảm giá. \\
\texttt{remaining\_quantity} / \texttt{expiry\_date} & & Giới hạn phát hành và hiệu lực. \\
\midrule
\multicolumn{3}{@{}l}{\cellcolor{tblaltbg}\textbf{Order} (Đơn hàng chính) --- FR5.2} \\
\texttt{id} & PK & \\
\texttt{user\_id} & FK $\to$ User.id & Người mua. \\
\texttt{total\_amount} / \texttt{created\_at} & & Giá trị và thời điểm đặt hàng. \\
\midrule
\multicolumn{3}{@{}l}{\cellcolor{tblaltbg}\textbf{SubOrder} (Đơn hàng con) --- tách theo gian hàng, FR5.2 và FR5.4} \\
\texttt{id} & PK & \\
\texttt{order\_id} & FK $\to$ Order.id & \\
\texttt{shop\_id} & FK $\to$ Shop.id & Mỗi đơn con thuộc đúng một gian hàng. \\
\texttt{status} & enum & Chờ xác nhận / đang giao / hoàn tất / hủy / hoàn tiền. \\
\texttt{shipping\_fee} / \texttt{total\_amount} & & Tính riêng cho từng đơn con (FR8.2). \\
\midrule
\multicolumn{3}{@{}l}{\cellcolor{tblaltbg}\textbf{OrderItem} (Chi tiết đơn hàng)} \\
\texttt{id} & PK & \\
\texttt{sub\_order\_id} & FK $\to$ SubOrder.id & \\
\texttt{batch\_id} & FK $\to$ Batch.id & Trừ tồn theo đúng lô, giữ liên kết truy xuất nguồn gốc (FR4.3). \\
\texttt{quantity} / \texttt{unit\_price} & & \\
\midrule
\multicolumn{3}{@{}l}{\cellcolor{tblaltbg}\textbf{Payment} (Thanh toán) --- dùng chung B2C/B2B, FR7.1--FR7.4} \\
\texttt{id} & PK & \\
\texttt{order\_id} & FK $\to$ Order.id, nullable & Luồng bán lẻ B2C. \\
\texttt{contract\_id} & FK $\to$ Contract.id, nullable & Luồng bán sỉ B2B. \\
\texttt{installment\_id} & FK $\to$ DeliveryInstallment.id, nullable & Thanh toán theo từng đợt giao. \\
\texttt{method} & enum & COD / chuyển khoản / ví điện tử. \\
\texttt{status} & enum & Trạng thái thanh toán. \\
\texttt{holding\_type} & enum & \emph{standard} hoặc \emph{escrow}. \\
\texttt{escrow\_status} & enum & Chưa ký quỹ / đang giữ / đã giải ngân / hoàn trả --- FR7.4. \\
\texttt{gateway\_transaction\_id} & & Mã đối soát với cổng thanh toán trung gian. \\
\texttt{paid\_at} & datetime & \\
\midrule
\multicolumn{3}{@{}l}{\cellcolor{tblaltbg}\textbf{Review} (Đánh giá) --- FR5.5} \\
\texttt{id} & PK & \\
\texttt{type} & enum & \emph{product} hoặc \emph{shop}. \\
\texttt{sub\_order\_id} & FK $\to$ SubOrder.id & Xác định gian hàng được đánh giá và bảo đảm điều kiện đã mua hàng. \\
\texttt{order\_item\_id} & FK $\to$ OrderItem.id, nullable & Chỉ dùng khi \texttt{type = product}. \\
\texttt{user\_id} & FK $\to$ User.id & Người đánh giá. \\
\texttt{rating} & 1--5 & Điểm đánh giá đơn lẻ. \\
\texttt{content} / \texttt{images} & & Nhận xét và ảnh kèm theo. \\
\texttt{moderation\_status} & enum & Công khai / chờ duyệt / đã ẩn --- đánh giá bị ẩn không tính vào điểm uy tín. \\
\texttt{seller\_reply} & nullable & Phản hồi công khai một lần của nhà bán. \\
\end{longtable}
\endgroup
Quan hệ: một Order phân rã thành nhiều SubOrder (1--n, mỗi
SubOrder gắn đúng một Shop); mỗi SubOrder có nhiều
OrderItem; Payment gắn với Order (1--1 hoặc 1--n nếu cho
thanh toán từng phần); Review luôn tham chiếu một SubOrder đã hoàn tất ---
kèm theo một dòng OrderItem khi đánh giá sản phẩm --- để đảm bảo chỉ
đánh giá sau khi đã mua (ràng buộc nghiệp vụ), mỗi đơn hàng con cho phép tối
đa một đánh giá gian hàng và một đánh giá cho mỗi sản phẩm đã mua.
 
\paragraph*{Nhóm 4 --- Giao dịch bán sỉ B2B (ứng với FR6)}

\noindent Bảng~\ref{tab:db-nhom4} trình bày đặc tả các thực thể thuộc nhóm này.
 
\begingroup
\small
\renewcommand{\arraystretch}{1.2}
\setlength{\tabcolsep}{5pt}
\begin{longtable}{L{4.9cm} L{3.8cm} L{6.0cm}}
\caption{Đặc tả thực thể nhóm 4: giao dịch bán sỉ B2B}
\label{tab:db-nhom4} \\
\toprule
\rowcolor{tblheadbg}
\textbf{Trường} & \textbf{Kiểu / Ràng buộc} & \textbf{Mô tả \& FR liên quan} \\
\midrule
\endfirsthead
\toprule
\rowcolor{tblheadbg}
\textbf{Trường} & \textbf{Kiểu / Ràng buộc} & \textbf{Mô tả \& FR liên quan} \\
\midrule
\endhead
\bottomrule
\endfoot
\multicolumn{3}{@{}l}{\cellcolor{tblaltbg}\textbf{RFQ} (Yêu cầu báo giá) --- FR6.1} \\
\texttt{id} & PK & \\
\texttt{business\_id} & FK $\to$ Business.id & Bên mua. \\
\texttt{shop\_id} & FK $\to$ Shop.id & Bên bán nhận yêu cầu. \\
\texttt{product\_id} & FK $\to$ Product.id & \\
\texttt{requested\_quantity} & & Dùng để đối chiếu MOQ của nhà bán. \\
\texttt{desired\_delivery\_time} & & \\
\texttt{quality\_standard} & & VietGAP/GlobalGAP hoặc yêu cầu riêng. \\
\texttt{status} & enum & Mới / đang đàm phán / chốt / từ chối / hết hạn. \\
\midrule
\multicolumn{3}{@{}l}{\cellcolor{tblaltbg}\textbf{RFQNegotiation} (Lịch sử đàm phán RFQ) --- lịch sử đàm phán, FR6.2} \\
\texttt{id} & PK & \\
\texttt{rfq\_id} & FK $\to$ RFQ.id & \\
\texttt{sender\_id} & FK $\to$ User.id & Bên gửi đề xuất. \\
\texttt{proposed\_price} / \texttt{proposed\_quantity} & & Nội dung thương lượng từng lượt. \\
\texttt{content} / \texttt{sent\_at} & & \\
\midrule
\multicolumn{3}{@{}l}{\cellcolor{tblaltbg}\textbf{Contract} (Hợp đồng) --- FR6.3} \\
\texttt{id} & PK & \\
\texttt{rfq\_id} & FK $\to$ RFQ.id & Một RFQ đã chốt sinh đúng một hợp đồng. \\
\texttt{agreed\_price} / \texttt{total\_quantity} & & Điều khoản thương mại đã thống nhất. \\
\texttt{payment\_terms} & enum & Net 15, Net 30... --- FR7.3. \\
\texttt{contract\_file} & & Bản hợp đồng điện tử đã sinh. \\
\texttt{signing\_status} / \texttt{signed\_at} & & Chờ ký / đã ký / hủy. \\
\midrule
\multicolumn{3}{@{}l}{\cellcolor{tblaltbg}\textbf{DeliveryInstallment} (Đợt giao hàng) --- FR6.5} \\
\texttt{id} & PK & \\
\texttt{contract\_id} & FK $\to$ Contract.id & \\
\texttt{quantity} & & Sản lượng của đợt giao. \\
\texttt{planned\_delivery\_date} & date & \\
\texttt{status} & enum & Chờ / đang giao / đã giao. \\
\midrule
\multicolumn{3}{@{}l}{\cellcolor{tblaltbg}\textbf{DebtLedger} (Sổ công nợ) --- FR6.5 và FR7.3} \\
\texttt{id} & PK & \\
\texttt{contract\_id} & FK $\to$ Contract.id & \\
\texttt{amount\_due} / \texttt{due\_date} & & Khoản phải trả theo kỳ hạn đã ký. \\
\texttt{status} & enum & Chưa đến hạn / quá hạn / đã tất toán. \\
\end{longtable}
\endgroup
Quan hệ: RFQ 1--n RFQNegotiation; một RFQ đã chốt sinh đúng một Contract (1--1);
Contract 1--n DeliveryInstallment; Contract 1--n DebtLedger (theo từng đợt hoặc gộp);
Contract/DeliveryInstallment 1--n Payment (Nhóm 3) --- bảng Payment dùng chung cho
cả hai luồng nhờ bộ khóa ngoại nullable, đồng thời lưu trạng thái ký quỹ
escrow của từng đợt giao theo FR7.4.
 
\paragraph*{Nhóm 5 --- Vận chuyển \& khiếu nại (ứng với FR8)}
 
\begin{itemize}
    \item \textbf{Shipment} (Vận chuyển): \texttt{id} (PK),
    \texttt{sub\_order\_id} (FK $\to$ SubOrder.id, nullable nếu gắn
    Contract/DeliveryInstallment cho B2B), \texttt{installment\_id}
    (FK $\to$ DeliveryInstallment.id, nullable), \texttt{carrier},
    \texttt{tracking\_number}, \texttt{status}.
    \item \textbf{ShipmentStatusLog} (Lịch sử vận chuyển): \texttt{id} (PK),
    \texttt{shipment\_id} (FK $\to$ Shipment.id), \texttt{status},
    \texttt{timestamp}, \texttt{note}.
    \item \textbf{Complaint} (Khiếu nại): \texttt{id} (PK),
    \texttt{sub\_order\_id} (FK $\to$ SubOrder.id, nullable),
    \texttt{contract\_id} (FK $\to$ Contract.id, nullable),
    \texttt{sender\_id} (FK $\to$ User.id), \texttt{content},
    \texttt{status} (mới/đang xử lý/đã xử lý), \texttt{handled\_by\_admin\_id}
    (FK $\to$ User.id, nullable), \texttt{resolution}.
\end{itemize}
Ghi chú: Shipment và Complaint dùng khóa ngoại \emph{nullable kép} (B2C qua
SubOrder hoặc B2B qua Contract/DeliveryInstallment) để tái sử dụng chung một bảng
cho cả hai luồng giao dịch, tránh trùng lặp cấu trúc.
 
\paragraph*{Nhóm 6 --- Tương tác, thông báo (ứng với FR9, FR11)}
 
\begin{itemize}
    \item \textbf{Conversation} (Cuộc trò chuyện): \texttt{id} (PK),
    \texttt{buyer\_id} (FK $\to$ User.id), \texttt{shop\_id}
    (FK $\to$ Shop.id).
    \item \textbf{Message} (Tin nhắn): \texttt{id} (PK),
    \texttt{conversation\_id} (FK $\to$ Conversation.id),
    \texttt{sender\_id} (FK $\to$ User.id), \texttt{content},
    \texttt{sent\_at}.
    \item \textbf{Wishlist} (Danh sách yêu thích): \texttt{user\_id}
    (FK $\to$ User.id), \texttt{product\_id} (FK $\to$ Product.id) ---
    khóa chính kép (composite PK).
    \item \textbf{ShopFollow} (Theo dõi gian hàng): \texttt{user\_id}
    (FK $\to$ User.id), \texttt{shop\_id} (FK $\to$ Shop.id)
    --- khóa chính kép.
    \item \textbf{Notification} (Thông báo): \texttt{id} (PK),
    \texttt{user\_id} (FK $\to$ User.id), \texttt{type} (đơn
    hàng/RFQ/chứng nhận/tồn kho/khuyến mãi...), \texttt{content},
    \texttt{is\_read}, \texttt{created\_at}.
\end{itemize}
 
\paragraph*{Nhóm 7 --- Hỗ trợ AI (ứng với FR12)}
 
\begin{itemize}
    \item \textbf{RecommendationLog} (Lịch sử gợi ý): \texttt{id} (PK),
    \texttt{user\_id} (FK $\to$ User.id), \texttt{product\_id}
    (FK $\to$ Product.id), \texttt{score}, \texttt{created\_at} ---
    phục vụ FR12.1, dùng làm dữ liệu đánh giá độ chính xác mô hình.
    \item \textbf{DemandPriceForecast} (Dự báo giá và nhu cầu): \texttt{id} (PK),
    \texttt{category\_id} (FK $\to$ Category.id), \texttt{forecast\_period}
    (tuần/tháng), \texttt{forecast\_price}, \texttt{forecast\_demand},
    \texttt{created\_at} --- phục vụ FR12.2.
    \item \textbf{AnomalyAlert} (Cảnh báo bất thường): \texttt{id} (PK),
    \texttt{target\_type} (đơn hàng/đánh giá/giao dịch...),
    \texttt{target\_id}, \texttt{description}, \texttt{severity} (thấp/trung
    bình/cao), \texttt{processing\_status}, \texttt{admin\_verdict}
    (đúng vi phạm/không phải vi phạm --- dữ liệu vòng phản hồi để đo tỷ lệ
    dương tính giả tại Bảng~\ref{tab:nguongAI}), \texttt{created\_at} --- phục
    vụ FR12.3, liên hệ FR10.2 (giám sát vi phạm).
\end{itemize}

\paragraph*{Nhóm 8 --- Quản trị, vi phạm \& lưu vết (ứng với FR10, NFR9.1)}

\begin{itemize}
    \item \textbf{AdminAuditLog} (Nhật ký quản trị): \texttt{id} (PK),
    \texttt{admin\_id} (FK $\to$ User.id), \texttt{action}
    (duyệt/từ chối/khóa/mở khóa/chỉnh sửa/giải ngân...),
    \texttt{target\_type} (nhà bán/sản phẩm/chứng nhận/đánh giá/khiếu
    nại...), \texttt{target\_id}, \texttt{reason}, \texttt{timestamp}
    --- lưu vết mọi thao tác quản trị theo FR10.1--FR10.3, đồng thời là căn
    cứ rà soát cho ràng buộc bất biến NFR9.1.
    \item \textbf{Violation} (Vi phạm): \texttt{id} (PK),
    \texttt{target\_id} (FK $\to$ User.id hoặc Shop.id),
    \texttt{violation\_type} (giả mạo chứng nhận/sai thông tin sản phẩm/thao
    túng đánh giá/gian lận đơn hàng...), \texttt{severity},
    \texttt{sanction} (nhắc nhở/cảnh cáo công khai/tạm khóa/khóa
    vĩnh viễn), \texttt{alert\_id} (FK $\to$ AnomalyAlert.id,
    nullable --- nếu vi phạm xuất phát từ cảnh báo AI),
    \texttt{handled\_by\_admin\_id} (FK $\to$ User.id),
    \texttt{handled\_at} --- hiện thực quy trình xử lý vi phạm phân cấp tại
    mục~\ref{sec:quytrinhnghiepvu} (Quản trị hệ thống) và FR10.2.
\end{itemize}
Quan hệ: một AnomalyAlert có thể dẫn tới không hoặc một Violation (0--1);
mỗi Violation và mỗi thao tác kiểm duyệt đều sinh ít nhất một bản ghi
AdminAuditLog, đảm bảo mọi quyết định quản trị đều truy vết được người thực
hiện và lý do.

\subsubsection*{Sơ đồ quan hệ thực thể (ERD)}
\label{subsec:erd}

Dựa trên các thực thể và mối quan hệ đã đặc tả, mô hình quan hệ thực thể (ERD)
tổng quát của hệ thống Farmery được thể hiện tại Hình~\ref{fig:erd}. Sơ đồ tổng
hợp mối liên kết giữa 8 nhóm thực thể nghiệp vụ: tài khoản và phân quyền người
dùng, quản lý gian hàng và truy xuất nguồn gốc, chu trình giao dịch bán lẻ B2C
và bán sỉ B2B, vận chuyển, khiếu nại, tương tác người dùng, các bảng phục vụ
tính năng AI, cùng nhóm bảng quản trị/lưu vết vi phạm.

\begin{figure}[H]
\centering
% \includegraphics[width=\textwidth]{image/erd.png}
\caption{Sơ đồ quan hệ thực thể (ERD) hệ thống Farmery}
\label{fig:erd}
\end{figure}
% TODO: nhóm bổ sung file ảnh sơ đồ ERD vào image/erd.png và bỏ chú thích dòng lệnh \includegraphics ở trên. Lưu ý vẽ theo bản đặc tả hiện hành tại mục~\ref{subsec:dacta-csdl} (8 nhóm thực thể, đã gồm Nhóm 8 quản trị/lưu vết, trường điểm uy tín của Shop và các khóa ngoại nullable của Payment cho luồng B2B).


\subsection{Thiết kế mô-đun AI}
\label{sec:thietkeAI}

\subsubsection*{Luồng dữ liệu tổng thể}

% TODO: nhóm bổ sung sơ đồ kiến trúc cụ thể (\includegraphics) minh họa luồng mô tả dưới đây, và chốt công nghệ/mô hình dùng cho từng bước khi hiện thực (mục~\ref{sec:hienthucAI}).
Cả bốn mô-đun AI (Chương~\ref{chap:coso}) đều đi theo một luồng dữ liệu chung gồm 5 bước:
\begin{enumerate}
    \item \textbf{Thu thập}: ghi nhận sự kiện từ các quy trình nghiệp vụ (lượt xem/mua sản phẩm, đơn hàng, đánh giá, hành vi hủy đơn) vào kho dữ liệu.
    \item \textbf{Tiền xử lý}: làm sạch, chuẩn hóa và trích xuất đặc trưng (feature engineering) theo từng bài toán.
    \item \textbf{Huấn luyện/suy luận}: huấn luyện định kỳ (batch, ví dụ theo tuần) đối với mô hình gợi ý/dự báo, hoặc suy luận theo thời gian thực (online inference) đối với phát hiện bất thường và chatbot.
    \item \textbf{Phục vụ kết quả}: trả kết quả qua API nội bộ để các mô-đun nghiệp vụ khác sử dụng (trang chủ, trang quản trị, trợ lý ảo).
    \item \textbf{Vòng phản hồi}: ghi nhận phản ứng của người dùng với kết quả AI (có nhấp vào gợi ý không, đánh giá đúng/sai của quản trị viên với cảnh báo gian lận) để làm dữ liệu huấn luyện lại, cải thiện mô hình theo thời gian.
\end{enumerate}

Danh sách dưới đây tóm tắt đầu vào/đầu ra riêng của từng mô-đun trong luồng chung này:
\begin{itemize}
    \item \textbf{Gợi ý sản phẩm}: đầu vào là lịch sử tương tác/mua hàng của người dùng và đặc trưng sản phẩm (danh mục, vùng trồng, mùa vụ); đầu ra là danh sách sản phẩm được xếp hạng theo mức độ phù hợp, hiển thị tại trang chủ và trang sản phẩm liên quan.
    \item \textbf{Dự báo nhu cầu/giá}: đầu vào là dữ liệu lịch sử đơn hàng, sản lượng theo mùa vụ và vùng trồng; đầu ra là ước lượng nhu cầu/giá cho giai đoạn tới, hiển thị cho nhà bán như gợi ý kế hoạch sản xuất và cho quản trị viên như dữ liệu cảnh báo dư cung/khan hiếm.
    \item \textbf{Phát hiện gian lận/bất thường}: đầu vào là đặc trưng hành vi giao dịch và đánh giá (tần suất, độ lệch giá, mẫu hủy đơn, mẫu đánh giá bất thường --- mục~\ref{subsec:danhgianghiepvu}); đầu ra là điểm bất thường và cảnh báo gửi tới quản trị viên để xem xét thủ công, liên hệ trực tiếp với quy trình quản trị hệ thống (mục~\ref{sec:quytrinhnghiepvu}).
\end{itemize}

\subsubsection*{Ngưỡng chất lượng mục tiêu}
\label{subsec:nguongAI}

Vì cả bốn mô-đun AI đều chưa có dữ liệu vận hành thực tế ở thời điểm thiết kế, mục này đặt ra các mục tiêu đo lường cụ thể --- thay vì chỉ mô tả định tính --- để làm căn cứ đối chiếu khi đánh giá kết quả thực nghiệm tại mục~\ref{sec:danhgiahieunang}.

\begingroup
\small
\renewcommand{\arraystretch}{1.3}
\setlength{\tabcolsep}{5pt}
\begin{longtable}{L{2.2cm} L{5.3cm} L{5.7cm}}
\caption{Ngưỡng chất lượng mục tiêu cho từng mô-đun AI}
\label{tab:nguongAI} \\
\toprule
\rowcolor{tblheadbg}
\textbf{Mô-đun} & \textbf{Chỉ số / ngưỡng mục tiêu} & \textbf{Căn cứ} \\
\midrule
\endfirsthead
\toprule
\rowcolor{tblheadbg}
\textbf{Mô-đun} & \textbf{Chỉ số / ngưỡng mục tiêu} & \textbf{Căn cứ} \\
\midrule
\endhead
\bottomrule
\endfoot
\textbf{Gợi ý sản phẩm} &
Tỷ lệ nhấp (CTR) vào danh sách gợi ý cá nhân hóa cao hơn tối thiểu 20\% so với danh sách đối chứng không cá nhân hóa (sản phẩm mới nhất/bán chạy chung), đo qua thử nghiệm A/B hoặc so sánh trước--sau khi bật tính năng. &
Đồ án không có đủ dữ liệu lịch sử lớn để đặt mục tiêu precision/recall tuyệt đối như các hệ thống gợi ý thương mại đã vận hành lâu năm; đo uplift CTR so với một baseline không cá nhân hóa khả thi hơn và vẫn phản ánh đúng giá trị thực tế mà tính năng cần tạo ra. \\
\rowcolor{tblaltbg}
\textbf{Dự báo nhu cầu/giá} &
Sai số phần trăm tuyệt đối trung bình (MAPE) dưới 20\% trên tập kiểm tra tách theo thời gian (time-based split), tương đương độ chính xác từ khoảng 80\% trở lên. &
Đặt thấp hơn một chút so với mức khoảng 82\% Farm2Market (2025) công bố cho bài toán dự báo giá nông sản tương tự \parencite{farm2market_2025}, do đây là công trình nghiên cứu học thuật có thể đã tối ưu trên bộ dữ liệu chọn lọc; 80\% là mục tiêu khả thi, kiểm chứng được trong phạm vi đồ án và vẫn đủ hữu ích để hỗ trợ ra quyết định. \\
\textbf{Phát hiện gian lận/bất thường} &
Tỷ lệ cảnh báo dương tính giả (false positive) trong số cảnh báo gửi tới quản trị viên không vượt quá 30\%, đo bằng tỷ lệ quản trị viên đánh dấu ``không phải vi phạm'' trên tổng số cảnh báo đã xử lý thủ công (vòng phản hồi ở mục trên). &
Đây là bài toán thiếu nhãn ``gian lận thật'' đầy đủ để tính precision/recall chuẩn ngay từ đầu, nên tỷ lệ false positive theo đánh giá thủ công là chỉ số khả thi nhất; ngưỡng 30\% được chấp nhận vì bỏ sót gian lận (false negative) bị đánh giá nghiêm trọng hơn chi phí xem xét thêm một số cảnh báo sai. \\
\rowcolor{tblaltbg}
\textbf{Trợ lý ảo (chatbot)} &
Ngưỡng độ tin cậy nhận diện ý định (intent confidence) tối thiểu 70\% để hệ thống tự trả lời; dưới ngưỡng này, hội thoại được chuyển tiếp (escalate) sang nhân sự CSKH thật (mô tả tại mục dưới). &
70\% là điểm cân bằng giữa hai rủi ro đối lập: ngưỡng quá thấp khiến chatbot trả lời sai ý định, gây mất lòng tin với nhóm nông dân/HTX vốn e ngại công nghệ (mục~\ref{sec:doituongnguoidung}); ngưỡng quá cao khiến phần lớn câu hỏi bị escalate không cần thiết, làm mất giá trị tự động hóa (NFR6.2). \\
\end{longtable}
\endgroup

\subsubsection*{Kiến trúc trợ lý ảo (chatbot)}

Dựa trên phân loại task-oriented/open-domain tại Chương~\ref{chap:coso}, trợ lý ảo Farmery được thiết kế theo kiến trúc task-oriented gồm ba thành phần chính:
\begin{itemize}
    \item \textbf{Nhận diện ý định (intent recognition)}: phân loại câu hỏi của người dùng vào một tập ý định giới hạn theo miền nghiệp vụ (ví dụ: tra cứu đơn hàng, hỏi truy xuất nguồn gốc, hướng dẫn đăng sản phẩm, hỏi chính sách đổi trả), trích xuất thực thể liên quan (mã đơn hàng, mã lô sản phẩm).
    \item \textbf{Cơ sở tri thức}: tổng hợp từ dữ liệu sản phẩm/chứng nhận/truy xuất nguồn gốc đang có trên nền tảng, nội dung Điều khoản dịch vụ (Phụ lục D) và các câu hỏi thường gặp theo từng vai trò người dùng (mục~\ref{sec:doituongnguoidung}).
    \item \textbf{Sinh phản hồi và chuyển tiếp}: sinh câu trả lời dựa trên ý định và cơ sở tri thức đã khớp; khi độ tin cậy nhận diện ý định thấp hoặc câu hỏi nằm ngoài tập ý định đã huấn luyện, hệ thống chuyển tiếp (escalate) sang nhân sự CSKH thật (mục~\ref{sec:luongtiepcan}) kèm theo lịch sử hội thoại để tránh người dùng phải lặp lại câu hỏi.
\end{itemize}
Với nhóm nông dân/HTX ít rành công nghệ (mục~\ref{sec:doituongnguoidung}), giao diện hội thoại ưu tiên câu trả lời ngắn, gợi ý sẵn các lựa chọn bấm thay vì bắt buộc gõ văn bản tự do, và hỗ trợ giọng nói/ngôn ngữ địa phương khi khả thi.

\subsection{Thiết kế API}
% TODO: danh sách endpoint chính theo nhóm FR (phương thức, đường dẫn, quyền truy cập RBAC, dữ liệu vào/ra); tài liệu hóa theo OpenAPI/Swagger (NFR5.2).

\subsection{Thiết kế UI/UX}
% TODO: wireframe/mockup các màn hình chính; lưu ý nguyên tắc tối giản cho nhóm nông dân/HTX (NFR6).

\subsection{Thiết kế test case}
\label{sec:kichbankiemthu}
% TODO: bảng test case theo từng nhóm FR (mã test, mô tả, dữ liệu vào, kết quả mong đợi); test case phi chức năng bám theo Bảng~\ref{tab:nguongNFR}.

\subsection{Các chính sách đối với xử lý của hệ thống và các bên thứ ba}
\label{subsec:chinhsachbenthuba}

\begin{itemize}
    \item \textbf{Cổng thanh toán trung gian:} mọi giao dịch thanh toán và ký quỹ escrow đi qua cổng thanh toán được cấp phép; hệ thống không lưu dữ liệu thẻ (NFR3.4), chỉ lưu mã giao dịch để đối soát và ra lệnh giải ngân khi đủ điều kiện (FR7.4).
    \item \textbf{Đơn vị vận chuyển thứ ba (3PL):} tích hợp qua REST API, mỗi lần gọi có thời gian chờ tối đa 5 giây và cơ chế thử lại khi lỗi tạm thời (NFR7.1); sự cố giao hàng được theo dõi trên bảng điều khiển giám sát (FR10.5).
    \item \textbf{Đăng nhập qua bên thứ ba:} hỗ trợ OAuth 2.0; tài khoản đăng nhập qua mạng xã hội vẫn phải xác thực OTP cho các thao tác nhạy cảm (FR1.1, FR1.2).
    \item \textbf{Xác minh chứng nhận:} giai đoạn MVP xác minh thủ công; khi có kết nối với cơ quan cấp chứng nhận sẽ chuyển sang đối chiếu tự động.
    \item \textbf{Xử lý nội bộ:} giới hạn tần suất thao tác (khóa 30 phút sau 5 lần sai), khóa tồn kho khi đặt hàng để chống bán vượt, lưu vết toàn bộ thao tác quản trị (NFR3.5, NFR1.2).
\end{itemize}
% TODO: bổ sung chính sách cụ thể theo điều khoản của từng đối tác sau khi chốt nhà cung cấp dịch vụ.

\section{Kết chương}

Chương này đã mô hình hóa bảy quy trình nghiệp vụ bằng lưu đồ, trình bày thiết kế cơ sở dữ liệu tám nhóm thực thể, thiết kế bốn mô-đun AI kèm ngưỡng chất lượng mục tiêu và các chính sách xử lý với bên thứ ba.
% TODO: cập nhật kết chương sau khi hoàn thiện use case, giải pháp công nghệ, kiến trúc, sitemap, lược đồ tuần tự/lớp, API, UI/UX và test case.
